%% ============================================================
%%  AttnFuse — Meeting Prep with Advisor
%%  Goal: explain the project in plain English, anticipate
%%  advisor questions, and lay out a path to publication.
%% ============================================================
\documentclass[11pt]{article}

\usepackage[margin=0.9in]{geometry}
\usepackage[T1]{fontenc}
\usepackage{lmodern}
\usepackage{microtype}
\usepackage{amsmath,amssymb}
\usepackage{graphicx}
\usepackage{booktabs}
\usepackage{listings}
\usepackage{xcolor}
\usepackage[colorlinks=true,linkcolor=blue!60!black,
            citecolor=blue!60!black,urlcolor=blue!60!black]{hyperref}
\usepackage{parskip}
\usepackage{enumitem}
\usepackage{titlesec}
\usepackage{tcolorbox}
\tcbuselibrary{breakable}

\titleformat{\section}{\Large\bfseries\color{blue!60!black}}{\thesection.}{0.6em}{}
\titleformat{\subsection}{\large\bfseries\color{black!80}}{\thesubsection.}{0.6em}{}

\definecolor{qbg}{HTML}{F2F4F8}
\definecolor{abg}{HTML}{FFF7E6}
\definecolor{tipbg}{HTML}{E8F4EA}

\newtcolorbox{question}[1]{%
  colback=qbg, colframe=blue!30!black, boxrule=0.4pt, arc=2pt,
  left=8pt, right=8pt, top=4pt, bottom=4pt,
  breakable, title=\textbf{Q: #1}, fonttitle=\bfseries\small}

\newtcolorbox{answer}{%
  colback=abg, colframe=orange!50!black, boxrule=0.4pt, arc=2pt,
  left=8pt, right=8pt, top=4pt, bottom=4pt,
  breakable, title=\textbf{Your answer}, fonttitle=\bfseries\small}

\newtcolorbox{layman}{%
  colback=tipbg, colframe=green!40!black, boxrule=0.4pt, arc=2pt,
  left=8pt, right=8pt, top=4pt, bottom=4pt,
  breakable, title=\textbf{In plain English}, fonttitle=\bfseries\small}

\lstset{
  language=Python,
  backgroundcolor=\color{black!4},
  basicstyle=\ttfamily\small,
  frame=single,
  rulecolor=\color{black!20},
  breaklines=true,
  showstringspaces=false,
  tabsize=2,
}

\title{\textbf{AttnFuse — Advisor Meeting Prep}\\[4pt]
       \large Layman walkthrough, anticipated questions, and path to publication}
\author{Varun Kumar Dasoju}
\date{}

\begin{document}
\maketitle

\noindent\textit{Read the green boxes for the simple version of every idea, the
gray-blue boxes for likely advisor questions, and the orange boxes for what to
say back. The last section sketches what would turn this from a class project
into a publishable paper.}

\tableofcontents
\bigskip

%% =============================================================
\section{The Project in Plain English}
%% =============================================================

\subsection{The one-paragraph pitch}

Modern AI language models spend most of their time inside a step called
\emph{attention}. Attention is the thing that lets the model decide which
other words it should ``look at'' when reading a given word. It is also the
slowest part: for a text of length $N$, attention naively builds an $N \times N$
table of scores, which is a memory disaster at long context. There is a
hand-written CUDA kernel called FlashAttention that solves this for the basic
case, but every time a researcher invents a new variant of attention
(sliding-window, ALiBi, RoPE, block sparse, etc.) they have to write a new
custom GPU kernel from scratch. \textbf{AttnFuse is a small programming
language embedded in Python that lets you describe the attention variant in
five lines of code; our compiler then automatically generates a fused GPU
kernel for it.} On an RTX~3090, the generated kernels are within
14\,--\,17\% of PyTorch's hand-tuned CUDA for the standard cases, and
17\,--\,38$\times$ faster than PyTorch for the cases PyTorch doesn't natively
accelerate.

\begin{layman}
Think of it like a recipe language for attention. The user writes a five-line
recipe; we generate the optimised GPU machine code that follows the recipe.
For common dishes (dense, causal) we are about as fast as the chef who has
been cooking them for years. For new dishes (sliding-window + ALiBi + RoPE +
custom bias) we are dozens of times faster than the only available
alternative, because that alternative is to throw the recipe out and cook
the whole meal by hand the slow way.
\end{layman}

\subsection{What's actually inside the box}

There are four pieces:

\begin{enumerate}[leftmargin=1.4em,itemsep=2pt]
  \item \textbf{A DSL frontend.} Seven Python ``combinators'' --- score, RoPE,
        causal mask, sliding window, ALiBi bias, additive bias, softmax ---
        that you snap together inside an \texttt{@af.attention}-decorated
        function. We trace the function symbolically (no GPU work yet) to
        capture a dataflow graph.
  \item \textbf{A two-level IR.} The high-level IR is a typed dataflow graph
        (TensorSym, ScoreOp, MaskOp, BiasOp, NormOp, MatMulPV). The low-level
        IR is a flat \texttt{TiledKernel} record with everything the
        codegenerator needs: mask kind, bias kind, RoPE flag, tile sizes,
        warp count, pipeline stages. The graph's SHA-1 signature is the
        kernel cache key.
  \item \textbf{A four-pass compiler.}
        \emph{Fuse} (recognise fusable patterns) $\to$
        \emph{Tile} (pick block sizes from an Ampere-tuned table) $\to$
        \emph{Lower} (collapse to TiledKernel) $\to$
        \emph{Codegen} (emit a single templated Triton kernel). Triton's JIT
        then specialises a distinct PTX binary per unique
        \texttt{tl.constexpr} combination, giving zero runtime branching.
  \item \textbf{A runtime.} Generated sources go to \texttt{\_generated/};
        the LaunchBundle precomputes dispatch metadata so per-call overhead is
        under one microsecond.
\end{enumerate}

\begin{layman}
We don't generate a new kernel for every new combination. We have one big
kernel template with switches baked in at compile time. When you ask for
``causal + ALiBi + RoPE,'' the Triton JIT bakes those flags in and the GPU
runs the specialised binary --- no \texttt{if}-statements in the hot loop.
\end{layman}

\subsection{The headline numbers}

Measured on RTX~3090, fp16, GPT-2 geometry, $N = 4096$:

\begin{itemize}[leftmargin=1.4em,itemsep=1pt]
  \item Dense: AttnFuse 59 TFLOPS vs SDPA 70 ($0.85\times$).
  \item Causal: AttnFuse 105 TFLOPS vs SDPA 122 ($0.86\times$).
  \item Sliding-window $W{=}256$: AttnFuse 298 TFLOPS vs SDPA 7.8 ($\mathbf{38\times}$).
  \item Causal + ALiBi: AttnFuse 98 TFLOPS vs SDPA 5.6 ($\mathbf{17\times}$).
  \item Fused RoPE: $1.18\times$ to $2.29\times$ faster than pre-processing RoPE.
  \item Peak memory at $N{=}4096$: 200~MB vs 3,176~MB naive ($\approx 32\times$ less).
  \item 25 GPU correctness tests passing; three dtypes (fp16, bf16, fp32).
\end{itemize}

\subsection{What's genuinely new}

\begin{enumerate}[leftmargin=1.4em,itemsep=2pt]
  \item \textbf{The DSL$+$compiler combination for attention.} Existing tools
        are either fixed (FlashAttention's two variants), too general
        (TVM, Halide), or new-this-year (PyTorch's \texttt{flex\_attention}).
        AttnFuse stakes out a different position --- a small dedicated IR
        with explicit introspection, ablation hooks, and per-combinator
        caching.
  \item \textbf{Fused RoPE inside the inner tile loop.} Rotate Q once before
        the loop; rotate K per-tile using a gather over rotate-half indices;
        all in registers, no HBM round-trip. To our knowledge, this is the
        first compiler-generated fused RoPE.
  \item \textbf{Variant-correct FLOP accounting and analytical roofline.}
        Most ``flash-attention TFLOPS'' numbers in the literature use the
        dense $4BHN^2D$ formula even for sub-quadratic kernels, producing
        tc\_pct values that exceed 100\%. We use $4BHN\min(2W,N)D$ for
        sliding-window, which is physically correct.
\end{enumerate}

%% =============================================================
\section{Anticipated Advisor Questions (with answers)}
%% =============================================================

\subsection{``What is the actual research contribution?''}

\begin{question}{Is this a compilers paper, a systems paper, or an engineering
project? What is the minimum publishable unit?}
\end{question}
\begin{answer}
Three candidates, strongest first:
(a)~the fused RoPE kernel and its codegen integration --- a focused systems
result with a clean before/after comparison, suitable for MLSys or a
workshop;
(b)~the combinator algebra plus two-level IR for attention --- a programming
languages / compilers angle, suitable for MAPL or a PLDI SRC paper;
(c)~the full DSL framework as an artifact paper for an MLSys-style track.

The cleanest first paper is (a) + a slimmer version of (c) as supporting
material: ``A composable IR for attention, with a novel fused-RoPE
specialisation.''
\end{answer}

\begin{question}{What's the delta vs xFormers, Liger Kernel, FlashAttention-2,
and PyTorch's new \texttt{flex\_attention}?}
\end{question}
\begin{answer}
xFormers and Liger ship hand-written kernels for a fixed set of variants ---
they don't expose a composable IR. FlashAttention-2 is exactly two variants
(dense, causal). The closest competitor is \texttt{flex\_attention} (PyTorch
2.5, 2024), which lets you supply a \texttt{score\_mod} Python function. Two
real differences: (i) \texttt{flex\_attention} cannot fuse pre-score
transformations like RoPE --- it only modifies the score, so RoPE still
needs preprocessing; (ii) AttnFuse exposes an explicit IR you can dump,
diff, and ablate, which we used to drive the tile-config study and the
roofline analysis. Different abstraction levels: \texttt{flex\_attention}
is a hook into PyTorch's compiler; AttnFuse is a small standalone DSL with
its own pipeline.
\end{answer}

\begin{question}{Has anyone done a compiler-generated attention DSL before?}
\end{question}
\begin{answer}
There is work in adjacent areas --- TVM's attention scheduling, Halide-style
auto-schedulers, MLIR's \texttt{linalg} dialect with transform-dialect
schedules, and FlexAttention as mentioned. But none of these combines a
combinator-based surface syntax + a typed dataflow IR + Triton codegen
specifically for attention, and none of them fuses RoPE. The novelty is
the combination, not any single component in isolation.
\end{answer}

\subsection{``Why Triton, why this IR design?''}

\begin{question}{Why generate Triton instead of writing your own MLIR
dialect, or writing CUDA directly?}
\end{question}
\begin{answer}
Triton is the right level of abstraction: high enough that we can template
a kernel and have its JIT specialise per constexpr, low enough to get
close-to-CUTLASS performance on Ampere. Writing an MLIR dialect adds
six months of infrastructure work for no expected performance win on the
RTX~3090 (the bottleneck is HBM, not codegen quality). Writing CUDA directly
would gain $\sim$15\% TFLOPS but cost the entire composability story --- we'd
end up shipping a fixed library of kernels rather than a DSL.
\end{answer}

\begin{question}{What does your IR actually buy you over just having a big
Python function with \texttt{if}-branches?}
\end{question}
\begin{answer}
Three things. (i)~Cache key: the graph signature is a stable hash, so two
users writing slightly different Python that produces the same graph share a
compiled kernel. (ii)~Validation: the fusion pass rejects unsupported
combinations (e.g.~softmax-of-non-score) at compile time, not at runtime.
(iii)~Ablation: we can run the same graph through multiple tile configs to
produce the heatmap in the report. None of these are possible with a
monolithic Python function.
\end{answer}

\subsection{``Are the speedups real?''}

\begin{question}{Your sliding-window speedup is vs SDPA's $\mathcal{O}(n^2)$
fallback. Isn't that an unfair comparison?}
\end{question}
\begin{answer}
Yes and no. Yes, in the sense that we are not beating a kernel that does the
same work --- we are beating PyTorch's only available execution path for
that variant. No, in the sense that this is exactly the user-visible
situation today: if you write \texttt{torch.nn.functional.scaled\_dot\_product\_attention}
with a non-trivial \texttt{attn\_mask}, you get the slow path. The fair
framing is ``AttnFuse delivers FlashAttention-class kernels for variants
PyTorch users currently cannot get accelerated at all,'' and that is what
we say in the paper.
\end{answer}

\begin{question}{Have you compared against hand-written FlashAttention for
sliding-window (e.g.~the Triton tutorial implementation)?}
\end{question}
\begin{answer}
We compare against \texttt{flash\_ref}, our own hand-written two-kernel
Triton FA2 reference, for dense and causal: AttnFuse is within 0.2\% of it.
For sliding-window there is no widely accepted ``hand-written FA2 sliding''
reference --- the closest is xFormers' \texttt{block\_diagonal\_attn}, which
we did not benchmark. This is a known limitation; adding xFormers to the
baseline matrix is on the to-do list for the paper.
\end{answer}

\begin{question}{RTX 3090 only --- what about A100, H100? The tile table is
hard-coded.}
\end{question}
\begin{answer}
True; the Ampere tile table will not give optimal numbers on Hopper
(sm\_90 has TMA, different SMEM budget, async copies). The IR is hardware-
neutral; what's hard-coded is just the lookup table in
\texttt{tiling.py}. Re-tuning it on H100 is the natural extension and is the
first item under ``Path to Publication.''
\end{answer}

\subsection{``What are you missing?''}

\begin{question}{No backward pass --- isn't this just inference?}
\end{question}
\begin{answer}
Correct. The backward kernel recomputes $m$ and $\ell$ from saved forward
statistics and has different tile scheduling (FA2 backward parallelises over
keys instead of queries). It's the largest missing piece; adding it would
roughly double the project's compiler-IR surface, but it's mechanical, not
research-novel. A workshop paper can defensibly cite ``forward only;
backward is straightforward future work,'' as the FA2 paper itself did
between v1 and v2.
\end{answer}

\begin{question}{No Nsight Compute? How do you defend the TFLOPS numbers?}
\end{question}
\begin{answer}
\texttt{ncu} was not installable on the evaluation host. We compute TFLOPS
analytically from latency $\times$ variant-correct FLOP count, and HBM
bandwidth from the analytical traffic (Q + K + V + O with $\min(2W,N)/N$
scaling for sliding-window). The HBM number is a lower bound that excludes
SMEM spills and pipeline overhead, which we explicitly note in the report.
For a venue submission, re-running on a box with \texttt{ncu} is a
day's work.
\end{answer}

\begin{question}{Why only one model geometry (GPT-2 small, 12 heads, $D{=}64$)?}
\end{question}
\begin{answer}
GPT-2-small and BERT-base were the proposal's evaluation targets. The kernel
parameterises on head\_dim ($32, 64, 96, 128, 256$) and num\_heads. For a
paper, the obvious extensions are LLaMA-style geometries ($D{=}128$,
GQA with groups$=$8) and Falcon-style ($D{=}64$, MQA). These are configuration
sweeps, not new code.
\end{answer}

\begin{question}{How big is your kernel-config search space? Are you really
``autotuning'' or just looking up a table?}
\end{question}
\begin{answer}
Honest answer: just a table. We did an offline ablation sweep
($\{BM, BN\} \times \text{warps} \times \text{stages}$) at $N{=}2048$ and
froze the best configs into \texttt{tiling.py}. We are not running
Triton's \texttt{autotune} at JIT time --- that would multiply compile
cost by 16$\times$. For a paper, the natural step is to integrate
\texttt{triton.autotune} per shape and report compile-time vs
runtime trade-off curves.
\end{answer}

\subsection{``How do you defend correctness?''}

\begin{question}{What's your test coverage? How do you know the fused RoPE
matches the unfused version?}
\end{question}
\begin{answer}
25 GPU unit tests cover every (variant $\times$ dtype) cell, plus dedicated
tests for the fused RoPE (compared against \texttt{naive\_attention} with
\texttt{apply\_rope} preprocessing; max absolute error $1.95 \times 10^{-3}$
at fp16, well within FA2's documented tolerance of $2 \times 10^{-2}$),
and for the sliding-window NaN guard. We do not yet have a property-based
fuzzer over combinator combinations; for a paper, that would be a
$\sim$200-line addition using Hypothesis.
\end{answer}

\begin{question}{What about numerical precision in low-precision dtypes? You
claim bf16 beats fp16 by 3\% --- is that signal or noise?}
\end{question}
\begin{answer}
Signal. The 3\% gap (108.4 vs 105.0 TFLOPS at $N=4096$ causal) is
reproducible across runs and consistent with the Triton/Ampere literature:
bf16 has a wider exponent range, which reduces masking-related underflow
in long-context softmax. We run 50 timed iterations after 5 warmups and
report the median; standard deviation across runs is under 1\%.
\end{answer}

%% =============================================================
\section{Hard / Critical Questions}
%% =============================================================

\begin{question}{If I strip away the DSL framing, what's left? Is the kernel
itself novel?}
\end{question}
\begin{answer}
Honest answer: the kernel-architecture is FlashAttention-2 with constexpr
specialisation --- not novel on its own. The fused RoPE codegen \emph{is}
novel. So the contribution is the combination: composable IR + codegen +
the RoPE specialisation. A reviewer who insists on a single-kernel novelty
result will not be impressed; a reviewer evaluating the IR + DSL + RoPE as
a system can be.
\end{answer}

\begin{question}{What stops someone from doing this in 200 lines of
\texttt{flex\_attention} score\_mod functions?}
\end{question}
\begin{answer}
The RoPE story is the strongest answer: \texttt{flex\_attention}'s
\texttt{score\_mod} runs \emph{on the score tensor}, after Q$\cdot$K$^\top$.
RoPE has to be applied \emph{to Q and K before} the matmul --- it's
mathematically not expressible as a score-mod function. So fused RoPE
specifically requires kernel-level support, which AttnFuse has and
\texttt{flex\_attention} does not.
\end{answer}

\begin{question}{What's the compile-time overhead in a realistic
training loop? 1.8 seconds per kernel sounds bad.}
\end{question}
\begin{answer}
1.8s is the cold-cache first call. Triton writes the PTX to
\texttt{\textasciitilde/.triton/cache/}, so the second program-launch is
near-instant. In a real training loop, you compile once per
(variant $\times$ shape) combination at the start; for a fixed-architecture
training run that's typically under 10 compilations total, amortised over
hundreds of millions of forward passes. The number to compare is
\textbf{compile-once-per-shape}, not per-call.
\end{answer}

\begin{question}{Are you actually going to win against PyTorch by the time
this gets reviewed? Their compiler stack moves fast.}
\end{question}
\begin{answer}
For dense/causal: no, and we don't claim to. For sliding-window + ALiBi +
RoPE compositions: PyTorch would need to either expand \texttt{flex\_attention}
to cover pre-score transformations, or add a new compile path for each
combination. Both are non-trivial. Our headline claim (``compiler-generated
fused RoPE'') is robust to PyTorch closing the gap on the score-mod path,
because RoPE is structurally not a score-mod.
\end{answer}

%% =============================================================
\section{Path to a Publishable Paper}
%% =============================================================

\subsection{Target venues, ranked by fit}

\begin{enumerate}[leftmargin=1.4em,itemsep=2pt]
  \item \textbf{MLSys (Mar deadline, Oct decisions).} Best fit: systems
        artifact paper with a compiler novelty and end-to-end measurements.
        Page limit allows the full IR + RoPE + evaluation story.
  \item \textbf{MAPL (PLDI workshop, Jun deadline).} Best fit for the
        ``combinator algebra + two-level IR + codegen'' angle if framed as
        a programming-languages contribution.
  \item \textbf{NeurIPS Workshop on Efficient ML / Systems for ML
        (Sep deadline).} Easier bar, good for a focused fused-RoPE result.
  \item \textbf{EuroMLSys / ML Systems track at PPoPP.} Solid backup.
\end{enumerate}

\subsection{Minimum additions to make the paper competitive}

Roughly in order of how much they strengthen the paper per unit of work:

\begin{enumerate}[leftmargin=1.4em,itemsep=2pt]
  \item \textbf{Add \texttt{flex\_attention} to the baseline matrix.}
        Direct head-to-head, including the RoPE case where it
        structurally cannot compete. Estimated effort: 1 week.
  \item \textbf{Re-tune the tile table on H100 (sm\_90) and report.}
        Demonstrates the IR is portable and not just a 3090 artifact.
        Need access to an H100. Effort: 1 week with hardware access.
  \item \textbf{Nsight Compute counter-based TFLOPS / HBM / occupancy.}
        Replaces our analytical estimates with measured numbers.
        Effort: 2 days on a box where \texttt{ncu} works.
  \item \textbf{Backward pass kernel + training-loop benchmark.}
        Closes the ``inference only'' criticism. This is the single
        largest piece of work: 3--4 weeks. Pays off in reviewer trust.
  \item \textbf{Block-sparse mask kind.} BigBird-style global+local+random,
        implemented as a new \texttt{MaskKind} with an index-buffer
        inner loop. Demonstrates the IR's extensibility. Effort: 2 weeks.
  \item \textbf{Larger model geometries (LLaMA $D{=}128$ GQA, Falcon MQA).}
        Configuration sweep, mostly free. Effort: 3 days.
  \item \textbf{Property-based correctness fuzzer.} Cheap reviewer
        comfort. Effort: 2 days.
  \item \textbf{End-to-end benchmark: full transformer forward, not just
        the attention op.} Demonstrates real-world relevance. Effort:
        1 week wiring through HuggingFace.
\end{enumerate}

\subsection{Suggested paper outline (8 pages)}

\begin{enumerate}[leftmargin=1.4em,itemsep=2pt]
  \item Introduction (1p): the gap, the contribution, headline numbers.
  \item Background (1p): online softmax, Triton, RoPE.
  \item AttnFuse design (2p): DSL, two-level IR, fusion + tiling +
        codegen passes, fused-RoPE codegen path with figure.
  \item Implementation (1p): tile-table derivation, dispatch
        optimisation, Triton pitfalls.
  \item Evaluation (2.5p): RTX 3090 + H100, all 5 variants, 3 dtypes,
        4 seq lengths, baselines = naive, SDPA, FA2 hand-ref,
        \texttt{flex\_attention}; roofline; ablation heatmap;
        fused-RoPE detailed breakdown; end-to-end transformer.
  \item Discussion + limitations + future work (0.5p).
\end{enumerate}

\subsection{Bold ``what would I do with another semester'' answer}

If asked for a vision pitch:
\begin{enumerate}[leftmargin=1.4em,itemsep=2pt]
  \item \textbf{Backward pass + training.} Turn AttnFuse into a training
        DSL, not just inference.
  \item \textbf{Hopper port (TMA + WGMMA + async pipeline).} Show the IR
        is hardware-neutral.
  \item \textbf{Quantised attention (fp8, int8) codegen.} The
        \texttt{ROPE\_KIND}-style constexpr trick generalises to
        quantisation modes with negligible compiler surface added.
  \item \textbf{Multi-GPU ring-attention codegen.} The IR can be lowered
        to a sequence-parallel variant; this is the biggest swing.
  \item \textbf{Integration with HuggingFace transformers via an
        \texttt{attn\_implementation="attnfuse"} backend.} Real users.
\end{enumerate}

%% =============================================================
\section{One-page Cheat Sheet (for the Meeting)}
%% =============================================================

\textbf{Pitch (30 seconds).}
``Attention is the slowest part of modern AI models. FlashAttention solves
it for two variants but requires hand-written CUDA. AttnFuse is a small DSL
in Python --- you write attention as five lines using combinators, we
compile to one fused Triton kernel. On RTX~3090, we match FA2 within 15\%
for dense and causal, and beat PyTorch by 17 to 38$\times$ for variants it
can't accelerate. Our most novel piece is the first compiler-generated
fused RoPE kernel.''

\textbf{If asked ``why a paper?''}
``Three angles: the IR design itself, the fused-RoPE codegen, and the
empirical demonstration that a DSL can match hand-tuned kernels for
common variants while opening up the long tail. The fused-RoPE result is
the strongest single new artifact; the rest is good supporting material.''

\textbf{If asked ``what's missing?''}
``Backward pass, H100 numbers, \texttt{flex\_attention} baseline, Nsight
Compute counters, and block-sparse. None of those are research-novel ---
they're systems-engineering work I can scope out in advance.''

\textbf{If asked ``what venue?''}
``MLSys is the best fit. MAPL workshop is a good rehearsal venue. NeurIPS
Efficient-ML workshop is a softer landing if the timing doesn't work.''

\textbf{Things to actively ask the advisor:}
\begin{enumerate}[leftmargin=1.4em,itemsep=1pt]
  \item Which of the three contribution framings ((a) RoPE-focused,
        (b)~IR-focused, (c) full system) does he see as the strongest?
  \item Does the department / group have H100 access I could use?
  \item Is he aware of any in-flight \texttt{flex\_attention} extensions
        that would moot our RoPE argument before submission?
  \item Would he be willing to be a co-author / advisor on submission?
        Confirming this early shapes scope.
\end{enumerate}

\end{document}
